\section{Signed Portfolios, Model Uncertainty, and Marked Losses}
\label{sec:risk_contract_extensions}
The nonnegative-output result does not cover every hedged portfolio. This section states a signed extension for a specified output location, an explicit failure for general portfolio weights, and sufficient uncertainty/noise contracts. The arguments use established block-power and quadratic-Lyapunov tools; they do not constitute a new general control method or authenticate a market loss calibration.

\begin{proposition}[Signed directed coupling with downstream-only output]
\label{prop:signed_target_coupling}
Let $A,B$ be fixed Schur-stable matrices, $L$ any real matrix, $\theta\ge0$, and
\[
J_\theta=\begin{bmatrix}A&0\\\theta L&B\end{bmatrix},\qquad C=[0,C_B],\quad C_B\ne0.
\]
Use the same fixed Euclidean initial-state ball and output units for every $\theta$. Then $M_C(J_\theta)=\sup_{t\ge0}\|CJ_\theta^t\|_2$ is nondecreasing in $\theta$, without sign restrictions on $A,B,L,C_B$. Consequently its exact initial-shock radius $B_{\rm loss}/M_C$ is nonincreasing.
\end{proposition}
\begin{proof}
For $t\ge1$, the lower-left block is $\theta S_t$, where $S_t=\sum_{j=0}^{t-1}B^{t-1-j}LA^j$. Thus $CJ_\theta^t=[\theta C_BS_t,C_BB^t]$. Its row Gram matrix is
\[
\theta^2 C_BS_tS_t^\top C_B^\top+C_BB^t(B^t)^\top C_B^\top.
\]
It increases in the positive-semidefinite order with $\theta$, so its largest eigenvalue and the spectral norm are nondecreasing. Taking the supremum over $t$, including the unchanged $t=0$ term, preserves the ordering.
\end{proof}
The output support is essential. For $\rho=0.85$, $h=2$, let
\[
J_\theta=\begin{bmatrix}\rho&h&0\\0&\rho&0\\0&\theta&\rho\end{bmatrix},\qquad C=[1,0,-1].
\]
All eigenvalues stay at $\rho$, while
\[
CJ_\theta^t=[\rho^t,(h-\theta)t\rho^{t-1},-\rho^t].
\]
At $\theta=h$ the middle term cancels and $M_C=\sqrt2$, whereas at $\theta=0$ the peak exceeds five. Increasing coupling can therefore enlarge a hedged portfolio's radius. This is an exact-output statement in fixed coordinates; it does not identify market hedging effectiveness from the observed VAR.

\subsection{Amplification Inside the Stated Multi-Asset Model}
An unrestricted Jordan example need not have a structural cohort parameterization. The stated leader--satellite equations supply a separate witness. Order states by asset as $(x_1,y_1,x_2,y_2)$, fix $a_i=0.5$, $\eta_i=0.2$, $q_i=0.2$, $K_{ii}=1$, $K_{12}=0$, and let $K_{21}=\theta$. Both diagonal blocks are
\[
A_i=\begin{bmatrix}0.5&0.2\\0.1&0.84\end{bmatrix},
\]
and the lower-left block is $\theta\left[\begin{smallmatrix}0&0.2\\0&0.04\end{smallmatrix}\right]$. The spectrum is independent of $\theta$ and is strictly inside the unit circle. For downstream price output $C=[0,0,1,0]$, a unit leader-memory shock produces loss $0.2\theta$ at one step. Hence $M_C\ge0.2\theta$ and the exact radius tends to zero as $\theta\to\infty$. This family satisfies the paper's linear cohort equations; economic no-arbitrage, feasible orders, impact calibration and actual portfolio exposures require separate justification.

For numerical certification, choose $h$ such that $q_h=\|J^h\|_2<1$. For $t=N+j+mh$, $0\le j<h$, submultiplicativity gives $\|CJ^t\|_2\le\|CJ^{N+j}\|_2q_h^m$. Therefore the infinite tail is bounded by $\max_{j<h}\|CJ^{N+j}\|_2$. If this bound is below the finite peak through $N-1$, that peak is certified. The examples use this tail check, rather than silently treating a finite horizon as infinite.

\begin{proposition}[A common-metric loss bound with changing coefficients]
\label{prop:robust_metric_loss}
Let $P\succ0$, $\|z\|_P=\|P^{1/2}z\|_2$, and
\[
z_t=J_tz_{t-1}+D_tw_t+e_t(z_{t-1}),\qquad y_t=C_tz_t.
\]
Suppose, throughout a region $\|z\|_P\le R$, that
\[
\|P^{1/2}J_tP^{-1/2}\|_2\le q_J,\quad
\|P^{1/2}D_t\|_2\le d,\quad
\|C_tP^{-1/2}\|_2\le c,\quad
\|e_t(z)\|_P\le\ell\|z\|_P.
\]
Let $q=q_J+\ell<1$, $\|w_t\|_2\le r$, $\|z_0\|_P\le r_0$, and $R\ge\max\{r_0,dr/(1-q)\}$. Then the region is invariant and
\[
\|y_t\|_2\le c\left(q^t r_0+\frac{dr(1-q^t)}{1-q}\right)
\le c\max\left\{r_0,\frac{dr}{1-q}\right\}.
\]
This is a sufficient uniform budget bound, not an exact optimal radius.
\end{proposition}
\begin{proof}
The declared operator and nonlinear bounds imply $\|z_t\|_P\le q\|z_{t-1}\|_P+dr$. The region condition gives $qR+dr\le R$, establishing invariance by induction. Iterating the scalar inequality gives the geometric expression. Multiplication by the output bound $c$ yields the result. The last expression is a convex combination of $r_0$ and $dr/(1-q)$.
\end{proof}
For a nominal stable $\widehat J$, solving $\widehat J^\top P\widehat J-P=-I$ supplies one common metric. If
$\|P^{1/2}(J_t-\widehat J)P^{-1/2}\|_2\le\varepsilon$, then $q_J\le\|P^{1/2}\widehat JP^{-1/2}\|_2+\varepsilon$. The uncertainty set, nonlinear remainder and exposures must actually satisfy these bounds; a fitted stable matrix does not establish them. Quadratic Lyapunov certificates are established control tools \citep{boyd1994lmi}.

Individual spectral stability cannot replace this common condition. The matrices $\left[\begin{smallmatrix}.5&2\\0&.5\end{smallmatrix}\right]$ and $\left[\begin{smallmatrix}.5&0\\2&.5\end{smallmatrix}\right]$ each have radius $.5$, while their alternating two-step product has radius above four. Coefficient switching can thus create growth even when every separately fitted matrix is stable. A weighted initial ball is also distinct from the earlier Euclidean ball: $\|z_0\|_P\le\sqrt{\lambda_{\max}(P)}\|z_0\|_2$. Coordinate changes require corresponding transformations of the metric, operators and uncertainty set.

\paragraph{Finite horizons without contraction.}
For a finite horizon, the same recursion admits any $q=q_J+\ell\ge0$:
\[
a_t=q^t r_0+dr\sum_{k=0}^{t-1}q^k,\qquad
R_H\ge\max_{0\le t\le H}a_t.
\]
A region satisfying this condition is invariant through $H$ by the same induction, and $\|y_t\|_2\le c a_t$ there. At $q=1$, $a_t=r_0+tdr$. Thus a declared finite-horizon contract can include a fundamental-price coordinate with a unit-root component; it supplies no indefinite contraction guarantee. Budget and nonlinear mark conditions must use this finite $R_H$.

\begin{proposition}[Finite-horizon Gaussian input contract]
\label{prop:finite_gaussian_risk}
For a declared horizon $H\ge1$, $0<\alpha<1$, and innovations with marginal law $w_t\sim\mathcal N(0,\sigma^2 I_m)$, define
\[
r_H=\sigma\sqrt{\chi^2_{m,1-\alpha/H}}.
\]
With probability at least $1-\alpha$, all $H$ input norms are at most $r_H$. If either the contracting or finite-horizon recursion's region and budget conditions hold with $r=r_H$, its bound holds simultaneously through $H$ on that event. Temporal independence is unnecessary for this union bound.
\end{proposition}
\begin{proof}
Each $\|w_t\|_2^2/\sigma^2$ has a chi-square law with $m$ degrees of freedom, so its tail beyond $r_H$ is $\alpha/H$. The union bound over $H$ steps gives failure probability at most $\alpha$. The deterministic recursion bound applies on the complementary event.
\end{proof}
The horizon and noise law are parts of the contract. If the required model/parameter bounds hold only on an event with failure probability at most $\beta$, a separate union bound gives combined coverage at least $1-\alpha-\beta$; the data must justify that additional coverage. Gaussian innovations do not satisfy an indefinitely bounded-input premise. Estimated noise scales, non-Gaussian tails, missed disturbances and model rejection require their own uncertainty analysis.

\begin{proposition}[Dollar mark mapping with a nonlinear remainder]
\label{prop:mark_to_dollar}
Let fixed positions $h_i$ hold instruments whose reference prices $P_i>0$ are denominated in the same budget currency. Let $b_i$ be actual log-price changes in basis points, with $|b_i|\le\Delta_i$. The marked profit change is
\[
\Pi(b)=\sum_i h_iP_i(e^{b_i/10000}-1),\qquad
\Pi_{\rm lin}(b)=\sum_i\frac{h_iP_i}{10000}b_i.
\]
Then
\[
|\Pi-\Pi_{\rm lin}|\le
\frac12\sum_i |h_i|P_i e^{\Delta_i/10000}(\Delta_i/10000)^2.
\]
If $b=Tz$, the dollar output row is $C=(h_iP_i/10000)_iT$. A linear budget plus this remainder and any separately bounded cash/execution costs must fit within the allocated capital budget.
\end{proposition}
\begin{proof}
Taylor's theorem gives $|e^v-1-v|\le e^{|v|}v^2/2$. Apply this componentwise and use the triangle inequality. Substitution of $b=Tz$ gives the stated dollar operator. Under $\|z\|_P\le R$, valid component bounds are $\Delta_i=R\|T_iP^{-1/2}\|_2$.
\end{proof}
Actual log-price displacement is distinct from latent pricing error relative to fundamentals. A stable error process can coexist with a moving fundamental price. If a price exactly tracks a Gaussian log-fundamental innovation $\xi$, then its pricing error stays zero, but a short position of size $|h|$ has marked loss $|h|P_0(e^\xi-1)$. For every finite budget $B$, the event $\xi>\log(1+B/(|h|P_0))$ has positive probability. Thus a deterministic mark-loss guarantee does not follow from zero pricing error. A finite sequence of doubling prices illustrates this omitted channel. Thus a loss operator on pricing error alone does not bound all mark risk. Fundamental-price and currency channels, changing positions, cash flows and margin rules must be included or separately bounded. The present proposition supplies a conditional unit/remainder mapping, rather than an authenticated institutional capital or execution guarantee.

\begin{table}[htbp]
\centering\small
\caption{Constructed fixed-coordinate examples, budget one in declared output units. The tail certificate bounds all times after the 150-step search; market calibration is not asserted.}\label{tab:risk_contract_examples}
\begin{adjustbox}{max width=\textwidth}
\begin{tabular}{lrrrrrr}\toprule
Family & Coupling & $\rho$ & Peak & $M_C$ & $1/M_C$ & Tail upper\\\midrule
Hedged & 0.0 & 0.850000 & 6 & 5.351112 & 0.186877 & 9.13e-09\\
Hedged & 1.0 & 0.850000 & 6 & 2.715136 & 0.368306 & 4.57e-09\\
Hedged & 2.0 & 0.850000 & 0 & 1.414214 & 0.707107 & 3.66e-11\\
Hedged & 3.0 & 0.850000 & 6 & 2.715136 & 0.368306 & 4.57e-09\\
Hedged & 4.0 & 0.850000 & 6 & 5.351112 & 0.186877 & 9.13e-09\\
Target market & 0.0 & 0.891133 & 0 & 1.000000 & 1.000000 & 1.45e-08\\
Target market & 2.0 & 0.891133 & 0 & 1.000000 & 1.000000 & 4.14e-07\\
Target market & 4.0 & 0.891133 & 4 & 1.262999 & 0.791766 & 8.27e-07\\
Target market & 8.0 & 0.891133 & 4 & 2.476201 & 0.403844 & 1.65e-06\\
Target market & 16.0 & 0.891133 & 4 & 4.927188 & 0.202956 & 3.31e-06\\
\bottomrule\end{tabular}
\end{adjustbox}
\end{table}
The structural radius in the table is computed from the unchanged diagonal blocks. Full-matrix floating-point eigenvalues are also retained digitally, because repeated eigenvalues can show small numerical discrepancies. Both families use fixed coordinates, the same initial-state ball within each family, and their stated output operators. These calculations illustrate the proved distinctions; they do not estimate real cross-impact or actual dealer safety.
